\documentclass[]{article}
\usepackage{lmodern}
\usepackage{amssymb,amsmath}
\usepackage{ifxetex,ifluatex}
\usepackage{fixltx2e} % provides \textsubscript
\ifnum 0\ifxetex 1\fi\ifluatex 1\fi=0 % if pdftex
  \usepackage[T1]{fontenc}
  \usepackage[utf8]{inputenc}
\else % if luatex or xelatex
  \ifxetex
    \usepackage{mathspec}
  \else
    \usepackage{fontspec}
  \fi
  \defaultfontfeatures{Ligatures=TeX,Scale=MatchLowercase}
\fi
% use upquote if available, for straight quotes in verbatim environments
\IfFileExists{upquote.sty}{\usepackage{upquote}}{}
% use microtype if available
\IfFileExists{microtype.sty}{%
\usepackage{microtype}
\UseMicrotypeSet[protrusion]{basicmath} % disable protrusion for tt fonts
}{}
\usepackage[left=1.1in, right=1.1in, top=1.1in, bottom=1.1in]{geometry}
\usepackage{hyperref}
\hypersetup{unicode=true,
            pdftitle={Stat 435 Lecture Notes 2a},
            pdfauthor={Xiongzhi Chen; Washington State University},
            pdfborder={0 0 0},
            breaklinks=true}
\urlstyle{same}  % don't use monospace font for urls
\usepackage{color}
\usepackage{fancyvrb}
\newcommand{\VerbBar}{|}
\newcommand{\VERB}{\Verb[commandchars=\\\{\}]}
\DefineVerbatimEnvironment{Highlighting}{Verbatim}{commandchars=\\\{\}}
% Add ',fontsize=\small' for more characters per line
\usepackage{framed}
\definecolor{shadecolor}{RGB}{248,248,248}
\newenvironment{Shaded}{\begin{snugshade}}{\end{snugshade}}
\newcommand{\KeywordTok}[1]{\textcolor[rgb]{0.13,0.29,0.53}{\textbf{{#1}}}}
\newcommand{\DataTypeTok}[1]{\textcolor[rgb]{0.13,0.29,0.53}{{#1}}}
\newcommand{\DecValTok}[1]{\textcolor[rgb]{0.00,0.00,0.81}{{#1}}}
\newcommand{\BaseNTok}[1]{\textcolor[rgb]{0.00,0.00,0.81}{{#1}}}
\newcommand{\FloatTok}[1]{\textcolor[rgb]{0.00,0.00,0.81}{{#1}}}
\newcommand{\ConstantTok}[1]{\textcolor[rgb]{0.00,0.00,0.00}{{#1}}}
\newcommand{\CharTok}[1]{\textcolor[rgb]{0.31,0.60,0.02}{{#1}}}
\newcommand{\SpecialCharTok}[1]{\textcolor[rgb]{0.00,0.00,0.00}{{#1}}}
\newcommand{\StringTok}[1]{\textcolor[rgb]{0.31,0.60,0.02}{{#1}}}
\newcommand{\VerbatimStringTok}[1]{\textcolor[rgb]{0.31,0.60,0.02}{{#1}}}
\newcommand{\SpecialStringTok}[1]{\textcolor[rgb]{0.31,0.60,0.02}{{#1}}}
\newcommand{\ImportTok}[1]{{#1}}
\newcommand{\CommentTok}[1]{\textcolor[rgb]{0.56,0.35,0.01}{\textit{{#1}}}}
\newcommand{\DocumentationTok}[1]{\textcolor[rgb]{0.56,0.35,0.01}{\textbf{\textit{{#1}}}}}
\newcommand{\AnnotationTok}[1]{\textcolor[rgb]{0.56,0.35,0.01}{\textbf{\textit{{#1}}}}}
\newcommand{\CommentVarTok}[1]{\textcolor[rgb]{0.56,0.35,0.01}{\textbf{\textit{{#1}}}}}
\newcommand{\OtherTok}[1]{\textcolor[rgb]{0.56,0.35,0.01}{{#1}}}
\newcommand{\FunctionTok}[1]{\textcolor[rgb]{0.00,0.00,0.00}{{#1}}}
\newcommand{\VariableTok}[1]{\textcolor[rgb]{0.00,0.00,0.00}{{#1}}}
\newcommand{\ControlFlowTok}[1]{\textcolor[rgb]{0.13,0.29,0.53}{\textbf{{#1}}}}
\newcommand{\OperatorTok}[1]{\textcolor[rgb]{0.81,0.36,0.00}{\textbf{{#1}}}}
\newcommand{\BuiltInTok}[1]{{#1}}
\newcommand{\ExtensionTok}[1]{{#1}}
\newcommand{\PreprocessorTok}[1]{\textcolor[rgb]{0.56,0.35,0.01}{\textit{{#1}}}}
\newcommand{\AttributeTok}[1]{\textcolor[rgb]{0.77,0.63,0.00}{{#1}}}
\newcommand{\RegionMarkerTok}[1]{{#1}}
\newcommand{\InformationTok}[1]{\textcolor[rgb]{0.56,0.35,0.01}{\textbf{\textit{{#1}}}}}
\newcommand{\WarningTok}[1]{\textcolor[rgb]{0.56,0.35,0.01}{\textbf{\textit{{#1}}}}}
\newcommand{\AlertTok}[1]{\textcolor[rgb]{0.94,0.16,0.16}{{#1}}}
\newcommand{\ErrorTok}[1]{\textcolor[rgb]{0.64,0.00,0.00}{\textbf{{#1}}}}
\newcommand{\NormalTok}[1]{{#1}}
\usepackage{graphicx,grffile}
\makeatletter
\def\maxwidth{\ifdim\Gin@nat@width>\linewidth\linewidth\else\Gin@nat@width\fi}
\def\maxheight{\ifdim\Gin@nat@height>\textheight\textheight\else\Gin@nat@height\fi}
\makeatother
% Scale images if necessary, so that they will not overflow the page
% margins by default, and it is still possible to overwrite the defaults
% using explicit options in \includegraphics[width, height, ...]{}
\setkeys{Gin}{width=\maxwidth,height=\maxheight,keepaspectratio}
\IfFileExists{parskip.sty}{%
\usepackage{parskip}
}{% else
\setlength{\parindent}{0pt}
\setlength{\parskip}{6pt plus 2pt minus 1pt}
}
\setlength{\emergencystretch}{3em}  % prevent overfull lines
\providecommand{\tightlist}{%
  \setlength{\itemsep}{0pt}\setlength{\parskip}{0pt}}
\setcounter{secnumdepth}{0}
% Redefines (sub)paragraphs to behave more like sections
\ifx\paragraph\undefined\else
\let\oldparagraph\paragraph
\renewcommand{\paragraph}[1]{\oldparagraph{#1}\mbox{}}
\fi
\ifx\subparagraph\undefined\else
\let\oldsubparagraph\subparagraph
\renewcommand{\subparagraph}[1]{\oldsubparagraph{#1}\mbox{}}
\fi

%%% Use protect on footnotes to avoid problems with footnotes in titles
\let\rmarkdownfootnote\footnote%
\def\footnote{\protect\rmarkdownfootnote}

%%% Change title format to be more compact
\usepackage{titling}

% Create subtitle command for use in maketitle
\newcommand{\subtitle}[1]{
  \posttitle{
    \begin{center}\large#1\end{center}
    }
}

\setlength{\droptitle}{-2em}

  \title{Stat 435 Lecture Notes 2a}
    \pretitle{\vspace{\droptitle}\centering\huge}
  \posttitle{\par}
    \author{Xiongzhi Chen \\ Washington State University}
    \preauthor{\centering\large\emph}
  \postauthor{\par}
    \date{}
    \predate{}\postdate{}
  
\usepackage{bbm}
\usepackage{amssymb}
\usepackage{amsmath}
\usepackage{graphicx,float}

\begin{document}
\maketitle

{
\setcounter{tocdepth}{2}
\tableofcontents
}
\section{}\label{section}

\section{Inference on coefficients}\label{inference-on-coefficients}

\subsection{Inference on estimates}\label{inference-on-estimates}

Since the observations \((x_i,y_i),i=1,\ldots,n\) for \((X,Y)\) are
random, so is the LSE \(\left(\hat{\beta}_0,\hat{\beta}_1\right)\).

\begin{itemize}
\item
  \(\left(\hat{\beta}_0,\hat{\beta}_1\right)\) is unbiased, i.e.,
  \(E(\hat{\beta}_0)=\beta_0\) and \(E(\hat{\beta}_1)=\beta_1\), when
  \(\varepsilon_i\)'s are uncorrelated
\item
  How accurate is \(\left(\hat{\beta}_0,\hat{\beta}_1\right)\) with
  respect to \(({\beta}_0,{\beta}_1)\)?
\end{itemize}

\subsection{Variability of estimates}\label{variability-of-estimates}

Recall
\[\hat{\beta}_1 = \frac{\sum (x_i - \bar{x})(y_i - \bar{y})}{\sum (x_i - \bar{x})^2}\]
and \(\hat{\beta}_0 = \bar{y} - \hat{\beta}_1 \bar{x}\).

When \(\varepsilon_i\)'s are uncorrelated,

\begin{itemize}
\tightlist
\item
  \([SE(\hat{\beta_0})]^2 = \sigma^2 \left[\frac{1}{n}+\frac{\bar{x}^2}{\sum_{i=1}^n(x_i - {\bar{x}})^2}\right]\)
\item
  \([SE(\hat{\beta_1})]^2 = \frac{\sigma^2}{\sum_{i=1}^n(x_i - \bar{x})^2}\)
\end{itemize}

How does the sample size \(n\) and the sample variance of \(x_i\)'s
affect these variances?

\subsection{Variability of estimates}\label{variability-of-estimates-1}

\begin{itemize}
\item
  Variances of \(\hat{\beta_0}\) and \(\hat{\beta_1}\) contain
  information on the accuracy of \(\hat{\beta}_0\) and \(\hat{\beta}_1\)
\item
  Without information on \(\sigma\), variances of \(\hat{\beta_0}\) and
  \(\hat{\beta_1}\) cannot be accurately assessed
\item
  An estimate of \(\sigma\) is
\end{itemize}

\[RSE = \sqrt{RSS/(n-2)} = \sqrt{\frac{\sum_{i=1}^n (y_i - \hat{y}_i)^2}{n-2}}\]

\subsection{CI from estimates}\label{ci-from-estimates}

The approximate \(95\%\) \emph{confidence interval (CI)}

\begin{itemize}
\tightlist
\item
  for \(\beta_0\) is \(\hat{\beta_0} \pm 2 \cdot SE(\hat{\beta_0})\)
\item
  for \(\beta_1\) is \(\hat{\beta_1} \pm 2 \cdot SE(\hat{\beta_1})\)
\end{itemize}

\emph{Note:} the above follows a general principle for constructing a CI
when the distribution of ``estimate minus parameter'' is symmetric
around \(0\)

\subsection{Testing the slope}\label{testing-the-slope}

\begin{itemize}
\tightlist
\item
  Recall the model \[Y = \beta_0 + \beta_1 X + \varepsilon\] with
  \(E(\varepsilon)=0\) and \(Var(\varepsilon)=\sigma^2\)
\item
  If \(\beta_1=0\), then \(Y = \beta_0 + \varepsilon\), and \(Y\) is
  independent of \(X\) (when \(\varepsilon\) is independent \(X\)).
\item
  Testing ``\(H_0: \beta_1=0\)'' often is equivalent to checking if
  ``\(H_0\): there is no relationship between \(X\) and \(Y\)''.
\end{itemize}

\emph{Caution:} The random error \(\varepsilon\) may not be independent
of \(X\) and \(Y\) due to latent dependence, which is common in genetics
studies.

\subsection{Testing the slope}\label{testing-the-slope-1}

\begin{itemize}
\tightlist
\item
  A test statistic for this purpose is the \emph{t-statistic}
  \[t = \frac{\hat{\beta}_1 - 0}{SE(\hat{\beta}_1)}\]
\item
  If \(SE(\hat{\beta}_1)\) is small, then relatively small
  \(\vert \hat{\beta}_1\vert\) provides strong evidence that
  \(\beta_1 \ne 0\)
\end{itemize}

\emph{Note:} When \(H_0: \beta_1=0\) holds, \(t\) approximately has a
\(t\)-distribution with \(n-2\) degrees of freedom when
\(\varepsilon_i\)'s are not much dependent on each other; when \(n\) is
large, a \(t\)-distribution will be close to a Gaussian distribution.

\subsection{Testing the slope}\label{testing-the-slope-2}

Model: \(E\)(\texttt{sales}) = \(\beta_0\) + \(\beta_1\) \texttt{TV}:

\begin{verbatim}
# A tibble: 2 x 5
  term        estimate std.error statistic  p.value
  <chr>          <dbl>     <dbl>     <dbl>    <dbl>
1 (Intercept)   7.03     0.458        15.4 1.41e-35
2 TV            0.0475   0.00269      17.7 1.47e-42
\end{verbatim}

\begin{itemize}
\tightlist
\item
  Is \(H_0: \beta_1=0\) retained or rejected? In either case, at which
  Type I error level?
\item
  How trustable are our conclusions?
\end{itemize}

\subsection{Testing the slope}\label{testing-the-slope-3}

Model \(E\)(\texttt{mpg}) = \(\beta_0\) + \(\beta_1\)
\texttt{horsepower}:

\begin{verbatim}
# A tibble: 2 x 5
  term        estimate std.error statistic   p.value
  <chr>          <dbl>     <dbl>     <dbl>     <dbl>
1 (Intercept)   39.9     0.717        55.7 1.22e-187
2 horsepower    -0.158   0.00645     -24.5 7.03e- 81
\end{verbatim}

\begin{itemize}
\tightlist
\item
  Is \(H_0: \beta_1=0\) retained or rejected? In either case, at which
  Type I error level?
\item
  How trustable are our conclusions?
\end{itemize}

\section{Model diagnostics}\label{model-diagnostics}

\subsection{Things to check}\label{things-to-check}

\begin{itemize}
\tightlist
\item
  Nonlinearity of relationship between response and predictor
\item
  Correlation of error terms
\item
  Non-constant variance of error terms
\item
  Outliers
\item
  High-leverage points
\end{itemize}

\subsection{Nonlinearity}\label{nonlinearity}

If the linear model \[Y=\beta_0 + \beta_1 X + \varepsilon\] is
plausible, then

\begin{itemize}
\tightlist
\item
  \(y_i = \beta_0 + \beta_1 x_i + \varepsilon_i\), where
  \(\varepsilon_i\) is a realization of \(\varepsilon\) for each pair
  \((x_i,y_i)\)
\item
  the residuals \(e_i\)'s\emph{, where
  \[e_i = y_i - \hat{y}_i, \quad \hat{y}_i=\hat{\beta}_0 + \hat{\beta}_1 x_i,\]
  }are random errors (as estimated realizations from \(\varepsilon\)),
  and should have no specific relationship with \(x_i\)'s
\end{itemize}

\emph{So, the residuals \(e_i\)'s should contain no specific pattern on
the fitted values \(\hat{y}_i\)'s or \(x_i\)'s if the model is
plausible.}

\subsection{Nonlinearity}\label{nonlinearity-1}

Model: \(E\)(\texttt{sales}) = \(\beta_0\) + \(\beta_1\) \texttt{TV}:

\begin{center}\includegraphics{Stat435LectureNotes2a_notes_files/figure-latex/unnamed-chunk-1-1} \end{center}

\subsection{Nonlinearity}\label{nonlinearity-2}

Model \(E\)(\texttt{mpg}) = \(\beta_0\) + \(\beta_1\)
\texttt{horsepower}:

\begin{center}\includegraphics{Stat435LectureNotes2a_notes_files/figure-latex/unnamed-chunk-2-1} \end{center}

\subsection{Check on error terms}\label{check-on-error-terms}

\begin{itemize}
\tightlist
\item
  Inference on the LSE \((\hat{\beta_0}, \hat{\beta_1})\) depends
  critically on \(SE(\hat{\beta_0})\) and \(SE(\hat{\beta_1})\), which
  depend on the unknown \(\sigma = \sqrt{Var(\varepsilon)}\).\\
\item
  Information on \(\varepsilon\), hence on \(\sigma\), is contained in
  the residuals \(e_i\) (as estimates of \(\varepsilon_i\))
\end{itemize}

Relatively accurate inference requires \(e_i\)'s to

\begin{itemize}
\tightlist
\item
  be \emph{uncorrelated}
\item
  have \emph{identical variance}
\end{itemize}

\emph{Otherwise, the formulae for \(SE(\hat{\beta_0})\) and
\(SE(\hat{\beta_1})\) are (usually) invalid, leading to invalid
inference.}

\subsection{Heterogeneous error
variances}\label{heterogeneous-error-variances}

Model: \(E\)(\texttt{sales}) = \(\beta_0\) + \(\beta_1\) \texttt{TV}:

\begin{center}\includegraphics{Stat435LectureNotes2a_notes_files/figure-latex/unnamed-chunk-3-1} \end{center}

\subsection{Heterogeneous error
variances}\label{heterogeneous-error-variances-1}

Model \(E\)(\texttt{mpg}) = \(\beta_0\) + \(\beta_1\)
\texttt{horsepower}:

\begin{center}\includegraphics{Stat435LectureNotes2a_notes_files/figure-latex/unnamed-chunk-4-1} \end{center}

\subsection{Outliers}\label{outliers}

\begin{itemize}
\tightlist
\item
  \emph{An outlier} is an observation \((x_i,y_i)\) for which \(y_i\) is
  far from the value predicted by the model
\item
  Observations whose studentized residuals are greater than \(3\) in
  absolute value are possible outliers
\item
  An outlier may significantly affect RSE (i.e., residual standard
  error) and \(R^2\)
\item
  Removing an outlier may or may not significantly affect the subsequent
  estimated regression line, and this is related to \emph{high-leverage
  points}
\end{itemize}

\subsection{Outliers}\label{outliers-1}

Model: \(E\)(\texttt{sales}) = \(\beta_0\) + \(\beta_1\) \texttt{TV}:

\begin{center}\includegraphics{Stat435LectureNotes2a_notes_files/figure-latex/unnamed-chunk-5-1} \end{center}

\subsection{Outliers}\label{outliers-2}

Model: \(E\)(\texttt{sales}) = \(\beta_0\) + \(\beta_1\) \texttt{TV}:

\begin{center}\includegraphics{Stat435LectureNotes2a_notes_files/figure-latex/unnamed-chunk-6-1} \end{center}

\subsection{Outliers}\label{outliers-3}

Model \(E\)(\texttt{mpg}) = \(\beta_0\) + \(\beta_1\)
\texttt{horsepower}:

\begin{center}\includegraphics{Stat435LectureNotes2a_notes_files/figure-latex/unnamed-chunk-7-1} \end{center}

\subsection{Outliers}\label{outliers-4}

Model \(E\)(\texttt{mpg}) = \(\beta_0\) + \(\beta_1\)
\texttt{horsepower}:

\begin{center}\includegraphics{Stat435LectureNotes2a_notes_files/figure-latex/unnamed-chunk-8-1} \end{center}

\subsection{High-leverage points}\label{high-leverage-points}

\begin{itemize}
\tightlist
\item
  \emph{A high-leverage point} is an observation \((x_i,y_i)\) for which
  \(x_i\) is unusual among all observations for \(X\)
\item
  Removing a high-leverage point often significantly affects the
  subsequent estimated regression line
\item
  Let \(p\) be the number of predictors in the model and \(n\) the
  sample size, if the leverage statistic for an observation greatly
  exceeds \((p+1)/n\), then it can be considered a high-leverage point
\end{itemize}

\subsection{High-leverage points}\label{high-leverage-points-1}

Model: \(E\)(\texttt{sales}) = \(\beta_0\) + \(\beta_1\) \texttt{TV}
(\(p=1\),\(n=397\), \(\tilde{h}=(p+1)/n \approx 0.005\)):
\vspace{-2.5cm}

\begin{center}\includegraphics{Stat435LectureNotes2a_notes_files/figure-latex/unnamed-chunk-9-1} \end{center}

\subsection{High-leverage points}\label{high-leverage-points-2}

Model \(E\)(\texttt{mpg}) = \(\beta_0\) + \(\beta_1\)
\texttt{horsepower} (\(p=1\),\(n=200\),\(\tilde{h}=(p+1)/n=0.01\)):
\vspace{-2.5cm}

\begin{center}\includegraphics{Stat435LectureNotes2a_notes_files/figure-latex/unnamed-chunk-10-1} \end{center}

\subsection{Q-Q Plot}\label{q-q-plot}

A \emph{Q-Q (quantile-quantile) plot} plots observed quantiles against
the quantiles of a theoretical distribution, and hence provides
information on whether observations under investigation have a
distribution that matches this theoretical distribution

\begin{itemize}
\tightlist
\item
  A normal Q-Q plot does this with the standard Normal distribution as
  the theoretical distribution
\item
  In a normal Q-Q plot, x-axis plots the theoretical quantiles from the
  standard Normal distribution with mean 0 and standard deviation 1
\end{itemize}

\subsection{Test on Normality}\label{test-on-normality}

\begin{itemize}
\tightlist
\item
  Model: \(E\)(\texttt{sales}) = \(\beta_0\) + \(\beta_1\) \texttt{TV}
\item
  Test Normality of random error
\end{itemize}

\begin{center}\includegraphics{Stat435LectureNotes2a_notes_files/figure-latex/unnamed-chunk-11-1} \end{center}

\subsection{Test on Normality}\label{test-on-normality-1}

\begin{itemize}
\tightlist
\item
  Model \(E\)(\texttt{mpg}) = \(\beta_0\) + \(\beta_1\)
  \texttt{horsepower}
\item
  Test Normality of random error
\end{itemize}

\begin{center}\includegraphics{Stat435LectureNotes2a_notes_files/figure-latex/unnamed-chunk-12-1} \end{center}

\subsection{Kolmogorov-Smirnov Test}\label{kolmogorov-smirnov-test}

\begin{itemize}
\tightlist
\item
  Model: \(E\)(\texttt{sales}) = \(\beta_0\) + \(\beta_1\) \texttt{TV}
\item
  \texttt{Fit1} is the object obtained from fitting the model
\item
  Test Normality of random error
\end{itemize}

\begin{verbatim}

    One-sample Kolmogorov-Smirnov test

data:  Fit1$residuals
D = 0.041533, p-value = 0.8806
alternative hypothesis: two-sided
\end{verbatim}

\subsection{Kolmogorov-Smirnov Test}\label{kolmogorov-smirnov-test-1}

\begin{itemize}
\tightlist
\item
  Model \(E\)(\texttt{mpg}) = \(\beta_0\) + \(\beta_1\)
  \texttt{horsepower}
\item
  \texttt{Fit2} is the object obtained from fitting the model
\item
  Test Normality of random error
\end{itemize}

\begin{verbatim}

    One-sample Kolmogorov-Smirnov test

data:  Fit2$residuals
D = 0.060525, p-value = 0.1131
alternative hypothesis: two-sided
\end{verbatim}

\subsection{Correlation of error
terms}\label{correlation-of-error-terms}

\emph{It is extremely important that the error terms are uncorrelated}.
Correlated error terms often present in time series data and in data
with latent variables. Such correlation affects

\begin{itemize}
\tightlist
\item
  testing if random errors are Normally distributed
\item
  variances of estimated coefficients and variance of random error
  term\\
\item
  Testing independence is a \emph{highly nontrival} issue in statistical
  learning
\end{itemize}

\subsection{Correlation of error
terms}\label{correlation-of-error-terms-1}

\begin{itemize}
\item
  The true model is \[Y = 1 + 2X + \varepsilon\] with \(n=1000\)
  observations \[y_{i} = 1+ 2 x_{i} + \varepsilon_{i}\]
\item
  The random errors are \emph{equally correlated}, such that
  \[\varepsilon_{i} =
  \sqrt{1-0.3}X_{i} + \sqrt{0.3} X_{0}\] and
  \(X_{0}, X_{1}, \ldots, X_{1000}\) are i.i.d. standard Normal.
\item
  Fit a simple linear model and obtain fitted values and residuals
\end{itemize}

\subsection{Correlation of error
terms}\label{correlation-of-error-terms-2}

\begin{center}\includegraphics{Stat435LectureNotes2a_notes_files/figure-latex/unnamed-chunk-15-1} \end{center}

\subsection{Correlation of error
terms}\label{correlation-of-error-terms-3}

For this example, when trying to check if the random errors are
independent or uncorrelated by a visual check, we see the following:

\begin{itemize}
\tightlist
\item
  no pattern in the left plot where each residuals is plotted against
  its index
\item
  the random errors are dependent, since if they were independent, the
  fitted values should be independent of the residuals
\item
  the random errors do not seem to be correlated with the fitted values
\end{itemize}

\section{Appendix}\label{appendix}

\subsection{True model}\label{true-model}

\begin{itemize}
\item
  For two quantitative random variables \(Y\) and \(X\), a simple linear
  model is \[E(Y) = \beta_0 + \beta_1 X,\] where \(\beta_0\)
  (\emph{intercept}) and \(\beta_1\) (\emph{slope}) are \emph{unknown,
  true model} parameters (or \emph{coefficients}), and \(\beta_1\) is
  called the \emph{regression coefficient}.
\item
  The above model is equivalent to
  \[Y = \beta_0 + \beta_1 X + \varepsilon \quad \text{ with } \quad E(\varepsilon)=0,Var(\varepsilon)=\sigma^2\]
  which is called the \emph{population regression line}.
\end{itemize}

\subsection{The least squares
estimate}\label{the-least-squares-estimate}

With observations \((x_i,y_i),i=1,\ldots,n\) for \((X,Y)\), the LS
method gives the \emph{least squares estimate (LSE)}:

\begin{itemize}
\item
  \(\hat{\beta}_1 = \frac{\sum (x_i - \bar{x})(y_i - \bar{y})}{\sum (x_i - \bar{x})^2}\)
  with \(\bar{y}=n^{-1}\sum_{i=1}^n y_i\)
\item
  \(\hat{\beta}_0 = \bar{y} - \hat{\beta}_1 \bar{x}\) with
  \(\bar{x}=n^{-1}\sum_{i=1}^n x_i\)
\end{itemize}

Namely, the fitted model is
\[\hat{y} = \hat{\beta}_0 + \hat{\beta}_1 x,\] also called the
\emph{least squares line}.

\subsection{License and session
Information}\label{license-and-session-information}

\href{http://math.wsu.edu/faculty/xchen/stat412/LICENSE.html}{License}

\begin{Shaded}
\begin{Highlighting}[]
\NormalTok{>}\StringTok{ }\KeywordTok{sessionInfo}\NormalTok{()}
\NormalTok{R version }\FloatTok{3.5.0} \NormalTok{(}\DecValTok{2018-04-23}\NormalTok{)}
\NormalTok{Platform:}\StringTok{ }\NormalTok{x86_64-w64-mingw32/}\KeywordTok{x64} \NormalTok{(}\DecValTok{64}\NormalTok{-bit)}
\NormalTok{Running under:}\StringTok{ }\NormalTok{Windows }\DecValTok{10} \KeywordTok{x64} \NormalTok{(build }\DecValTok{19045}\NormalTok{)}

\NormalTok{Matrix products:}\StringTok{ }\NormalTok{default}

\NormalTok{locale:}
\NormalTok{[}\DecValTok{1}\NormalTok{] LC_COLLATE=English_United States}\FloatTok{.1252} 
\NormalTok{[}\DecValTok{2}\NormalTok{] LC_CTYPE=English_United States}\FloatTok{.1252}   
\NormalTok{[}\DecValTok{3}\NormalTok{] LC_MONETARY=English_United States}\FloatTok{.1252}
\NormalTok{[}\DecValTok{4}\NormalTok{] LC_NUMERIC=C                          }
\NormalTok{[}\DecValTok{5}\NormalTok{] LC_TIME=English_United States}\FloatTok{.1252}    

\NormalTok{attached base packages:}
\NormalTok{[}\DecValTok{1}\NormalTok{] stats     graphics  grDevices utils     datasets  methods  }
\NormalTok{[}\DecValTok{7}\NormalTok{] base     }

\NormalTok{other attached packages:}
\NormalTok{[}\DecValTok{1}\NormalTok{] ggplot2_3}\FloatTok{.1.0} \NormalTok{broom_0}\FloatTok{.5.1}   \NormalTok{knitr_1}\FloatTok{.21}   

\NormalTok{loaded via a }\KeywordTok{namespace} \NormalTok{(and not attached):}
\StringTok{ }\NormalTok{[}\DecValTok{1}\NormalTok{] Rcpp_1}\FloatTok{.0.12}      \NormalTok{plyr_1}\FloatTok{.8.4}       \NormalTok{pillar_1}\FloatTok{.3.1}    
 \NormalTok{[}\DecValTok{4}\NormalTok{] compiler_3}\FloatTok{.5.0}   \NormalTok{tools_3}\FloatTok{.5.0}      \NormalTok{digest_0}\FloatTok{.6.18}   
 \NormalTok{[}\DecValTok{7}\NormalTok{] evaluate_0}\FloatTok{.13}    \NormalTok{tibble_2}\FloatTok{.1.3}     \NormalTok{nlme_3}\FloatTok{.1}\DecValTok{-137}    
\NormalTok{[}\DecValTok{10}\NormalTok{] gtable_0}\FloatTok{.2.0}     \NormalTok{lattice_0}\FloatTok{.20}\DecValTok{-35}  \NormalTok{pkgconfig_2}\FloatTok{.0.2} 
\NormalTok{[}\DecValTok{13}\NormalTok{] rlang_0}\FloatTok{.4.4}      \NormalTok{cli_1}\FloatTok{.0.1}        \NormalTok{rstudioapi_0}\FloatTok{.8}  
\NormalTok{[}\DecValTok{16}\NormalTok{] yaml_2}\FloatTok{.2.0}       \NormalTok{xfun_0}\FloatTok{.4}         \NormalTok{withr_2}\FloatTok{.1.2}     
\NormalTok{[}\DecValTok{19}\NormalTok{] dplyr_0}\FloatTok{.8.4}      \NormalTok{stringr_1}\FloatTok{.3.1}    \NormalTok{generics_0}\FloatTok{.0.2}  
\NormalTok{[}\DecValTok{22}\NormalTok{] grid_3}\FloatTok{.5.0}       \NormalTok{tidyselect_0}\FloatTok{.2.5} \NormalTok{glue_1}\FloatTok{.3.0}      
\NormalTok{[}\DecValTok{25}\NormalTok{] R6_2}\FloatTok{.3.0}         \NormalTok{fansi_0}\FloatTok{.4.0}      \NormalTok{rmarkdown_1}\FloatTok{.11}  
\NormalTok{[}\DecValTok{28}\NormalTok{] purrr_0}\FloatTok{.2.5}      \NormalTok{tidyr_0}\FloatTok{.8.2}      \NormalTok{magrittr_1}\FloatTok{.5}    
\NormalTok{[}\DecValTok{31}\NormalTok{] backports_1}\FloatTok{.1.3}  \NormalTok{scales_1}\FloatTok{.0.0}     \NormalTok{htmltools_0}\FloatTok{.3.6} 
\NormalTok{[}\DecValTok{34}\NormalTok{] assertthat_0}\FloatTok{.2.0} \NormalTok{colorspace_1}\FloatTok{.3}\DecValTok{-2} \NormalTok{labeling_0}\FloatTok{.3}    
\NormalTok{[}\DecValTok{37}\NormalTok{] utf8_1}\FloatTok{.1.4}       \NormalTok{stringi_1}\FloatTok{.2.4}    \NormalTok{lazyeval_0}\FloatTok{.2.1}  
\NormalTok{[}\DecValTok{40}\NormalTok{] munsell_0}\FloatTok{.5.0}    \NormalTok{crayon_1}\FloatTok{.3.4}    
\end{Highlighting}
\end{Shaded}


\end{document}
